\begin{claim}{C.27}{Proved}{An order-natural contraction of line centralizer loops}
Give $\operatorname{Homeo}(\R)$ its compact-open topology, let
$H\leq\operatorname{Homeo}(\R)$ be an arbitrary subgroup, and put
\[
              C=C_{\operatorname{Homeo}^+(\R)}(H).
\]
Every based loop in $C$ has an explicit based null-homotopy in $C$.
Orientation-reversing elements of $H$ are allowed. The construction
commutes with conjugation by every homeomorphism of the line and
requires no finiteness, differentiability, or local compactness of $H$.
\end{claim}
\begin{proof}
Let $t\mapsto f_t$ be a compact-open continuous loop in $C$, with
$f_0=f_1=\id$. For $0\leq r\leq1$ define
\[
 m_r(x)=\min_{0\leq u\leq r}f_u(x),\qquad
 M_r(x)=\max_{0\leq u\leq r}f_u(x).                    \tag{C.27.1}
\]
Both are increasing homeomorphisms. Indeed, for $x<y$, compactness
gives
\[
 \min_{0\leq u\leq1}\bigl(f_u(y)-f_u(x)\bigr)>0,
\]
which proves strict increase for each minimum and maximum. Their
properness is uniform in $r$: compactness of the inverse family
$\{f_u^{-1}:0\leq u\leq1\}$ gives, for every bounded target interval,
a common bound on all its inverse images. Continuity of (C.27.1),
jointly in $r$ and $x$, follows from uniform continuity on compact
parameter rectangles. Proper continuous strictly increasing maps of
$\R$ are homeomorphisms.

Set
\[
 F_{s,t}(x)=\operatorname{med}
      \bigl(m_{1-s}(x),f_t(x),M_{1-s}(x)\bigr)
   =\max\{m_{1-s}(x),\min\{f_t(x),M_{1-s}(x)\}\}.
                                                        \tag{C.27.2}
\]
Finite minima and maxima of increasing homeomorphisms are again
increasing homeomorphisms, by the same strict-increase and properness
arguments. Formula (C.27.2) is continuous in $(s,t,x)$, uniformly on
compact $x$-sets; it is therefore a continuous family in the
compact-open homeomorphism group. At $s=0$, the prefix interval is
the entire loop, so $m_1\leq f_t\leq M_1$ and $F_{0,t}=f_t$.
At $s=1$, $m_0=M_0=\id$, giving $F_{1,t}=\id$. Since $u=0$ occurs
in every prefix interval, $m_r\leq\id\leq M_r$, and hence
$F_{s,0}=F_{s,1}=\id$ for every $s$.

It remains to check the exact equivariance. An increasing
$h\in H$ fixes both envelopes under conjugation, because conjugation
commutes with their pointwise minimum and maximum and fixes every
$f_u$. A decreasing $h\in H$ interchanges $m_r$ and $M_r$ instead.
Every order-preserving or order-reversing bijection preserves the
median of three points. Thus both cases give
\[
                   h F_{s,t}h^{-1}=F_{s,t}
       \qquad(h\in H, 0\leq s,t\leq1).
\]
This proves that the entire based contraction lies in $C$. Replacing
the loop by its conjugate by an arbitrary line homeomorphism gives
the conjugate of the same formula, proving the asserted naturality.
\end{proof}

The surface problem in C.12 still requires a continuous construction
across different orbit strata. Applying (C.27.2) separately after
choosing coordinates on individual line orbits does not establish
continuity on the surface. No such gluing statement follows from
C.27, and none is used here.
